\section{Configuração Experimental e Estatísticas}
\label{sec:experimental}

\subsection{Amostragem e Avaliação}
Os métodos clássicos apresentam custo computacional heterogêneo e, por esse
motivo, a avaliação multi-sensor é conduzida em um subconjunto balanceado do
conjunto de dados, limitado a um número máximo de recortes por sensor e por
semente experimental. Esse desenho preserva comparabilidade entre sensores ao
mesmo tempo em que mantém a execução viável.

Para aprendizado profundo, a avaliação é restrita ao sensor Sentinel-2,
utilizando a partição de teste indicada na \Cref{tab:dataset-stats}. Quando a
comparação entre métodos clássicos e modelos de \gls{DL} é apresentada, os
resultados clássicos são filtrados para Sentinel-2, evitando confundir
diferenças metodológicas com diferenças de sensor.

\subsection{Métricas}
Reportamos \gls{PSNR}, \gls{SSIM}, \gls{RMSE}, \gls{SAM} e \gls{ERGAS} apenas em
pixels de lacuna. O \gls{RMSE} por banda espectral também é computado para
diagnóstico.

\subsection{Análise Estatística}
Adota-se correlação de Spearman para quantificar associações monotônicas entre
entropia local e métricas de qualidade, com correção por \gls{FDR} em múltiplos
testes. Para comparar distribuições de desempenho entre métodos, utiliza-se o
teste de Kruskal-Wallis, seguido de comparações par-a-par com correção
conservadora para múltiplas hipóteses e medidas de tamanho de efeito.

Como análise complementar, ajusta-se regressão robusta (\gls{RLM} com norma de
Huber) para modelar métricas como função das entropias em múltiplas escalas,
controlando por método e nível de ruído. Intervalos de confiança são obtidos
via bootstrap, com reamostragem em cluster por \texttt{patch\_id}, para reduzir
dependência intra-recorte.

\subsection{Ambiente de Execução}
Os experimentos foram executados em Rocky Linux 9.5 com NVIDIA A100 80GB PCIe,
driver 550.90.07 e CUDA 12.4/12.6.2.
